% !TeX root = ../main.tex
\subsection{تمرین \secnum{۵.۱}، تخمین توان سوئیچینگ تراشه}
\label{sec:assignment-power-chip}

\textbf{خواسته.}
تراشه‌ای از منطق تصادفی با \lr{Activity Factor} متوسط \(\alpha=0.1\)
در نظر بگیرید. ظرفیت سوئیچینگ بر واحد سطح در فرایند \lr{Standard Cell}
برابر \(c_A=\qty[per-mode=symbol]{450}{\pico\farad\per\square\milli\metre}\)
است. با مساحت \(A=\qty{70}{\square\milli\metre}\)، فرکانس
\(f=\qty{450}{\mega\hertz}\) و تغذیه \(V_{DD}=\qty{0.9}{\volt}\)،
توان دینامیکی تراشه را تخمین بزنید \cite[تمرین \secnum{۵.۱}، صفحه ۲۰۹]{weste}.

\textbf{حل.}
چگالی ظرفیت را در مساحت ضرب می‌کنیم تا ظرفیت کلِ پیش از وزن‌دهی
با فعالیت، یعنی \(C\)، به دست آید. واحد مساحت در ضرب ساده می‌شود.
\[
  C=c_AA=\qty{31500}{\pico\farad}=\qty{31.5}{\nano\farad}.
\]
برای اعمال فعالیت متوسط بر کل ظرفیت، مقدار \(0.1\) را نماینده
میانگین وزن‌دار \(\alpha=\sum_i\alpha_i C_i/\sum_i C_i\) می‌گیریم
که \(C_i\) و \(\alpha_i\) ظرفیت و فعالیت گره \(i\) هستند.
پس ظرفیت مؤثر \(\alpha C\) است. با فرض ناچیز بودن جریان
مستقیمِ گذرا، توان دینامیکی را با توان سوئیچینگ برآورد می‌کنیم
\cite[رابطه \secnum{۵.۱۰}، صفحه ۱۸۴]{weste}.
\[
  \begin{aligned}
    P_{\mathrm{dyn}}&\approx \alpha C V_{DD}^{2} f\\
      &=0.1(\qty{31.5e-9}{\farad})(\qty{0.9}{\volt})^{2}
             (\qty{450e6}{\hertz})\approx\qty{1.15}{\watt}.
  \end{aligned}
\]
حاصل ضرب فاراد و مربع ولت، ژول است. ضرب در هرتز آن را به وات
تبدیل می‌کند. مقدار نمایش‌داده‌شده به دو رقم اعشار گرد شده و
برآورد در مرتبه‌ی یک وات است.
